\subsection{平坦性判别准则的陈述}

设 A 是一个环，\(\mathfrak{I}\) 是 A 的双边理想，M 是左模，gr(A) 和 gr(M) 分别是由环 A 及模 M 配备 \(\mathfrak{I}\)-进滤过所相伴的分次环及分次 gr(A)-模（§ 2，第 3 节）。我们已经看到（同上），对于每个整数 \(n \geqslant 0\)，存在满射 Z-模同态

\[
\gamma_ {n} \colon (\mathfrak {I} ^ {n} / \mathfrak {I} ^ {n + 1}) \otimes_ {\mathrm{A} / \mathfrak {I}} (\mathrm{M} / \mathfrak {I} \mathrm{M}) \rightarrow \mathfrak {I} ^ {n} \mathrm{M} / \mathfrak {I} ^ {n + 1} \mathrm{M}
\]

以及一个次数为 0 的分次 gr(A)-模的分次同态

\[
\gamma_ {M} \colon \operatorname{gr} (A) \otimes_ {\mathbf {g r} _ {0} (A)} \operatorname{gr} _ {0} (M) \to \operatorname{gr} (M)
\]

其在 \(\mathrm{gr}_n(\mathbf{A})\otimes_{\mathbf{gr}_0(\mathbf{A})}\mathrm{gr}_0(\mathbf{M})\) 上的限制是 \(\gamma_{n}\) 对所有 \(n\) 成立，因而它是满射。

定理 1. 设 A 是交换环，\(\Im\) 是 A 的理想，M 是 A-模。考虑以下性质：

\begin{enumerate}
\def\labelenumi{(\roman{enumi})}
\item
  M 是平坦 A-模。
\item
  都有 \(\operatorname{Tor}_1^{\mathbf{A}}(\mathbf{N},\mathbf{M}) = 0\)，其中 \(\mathbf{N}\) 是被 \(\Im\) 消灭的 A-模。
\item
  M/3M 是平坦的 (A/3)-模，且典范映射 \(\Im \otimes_{\mathbf{A}} \mathbf{M} \to \Im \mathbf{M}\) 是双射（后一个条件等价于 \(\operatorname{Tor}_{1}^{\mathbf{A}}(\mathbf{A}/\Im, \mathbf{M}) = 0\)，这是由关系 \(\operatorname{Tor}_{1}^{\mathbf{A}}(\mathbf{A}, \mathbf{M}) = 0\) 以及正合列

\[
\operatorname{Tor} _ {1} ^ {\mathbf {A}} (\mathrm{A}, \mathrm{M}) \rightarrow \operatorname{Tor} _ {1} ^ {\mathbf {A}} (\mathrm{A} / \Im , \mathrm{M}) \rightarrow \Im \otimes_ {\mathrm{A}} \mathrm{M} \rightarrow \mathrm{M}).
\]

\item
  M/3M 是平坦的 (A/3)-模，且典范同态

\[
\gamma_ {\mathbf {M}} \colon \operatorname{gr} (\mathbf {A}) \otimes_ {\mathbf {g r} _ {0} (\mathbf {A})} \operatorname{gr} _ {0} (\mathbf {M}) \rightarrow \operatorname{gr} (\mathbf {M})
\]

是双射（§ 2，第 8 节的性质 (GR)）。

\item
  对所有 \(n \geqslant 1\)，\(\mathbf{M} / \Im^n\mathbf{M}\) 都是平坦的 \((\mathbf{A} / \Im^n)\)-模。
\end{enumerate}

\[
(i) \Rightarrow (i i) \Leftrightarrow (i i i) \Rightarrow (i v) \Leftrightarrow (v).
\]

此外，若 \(\mathfrak{I}\) 是幂零的，或者 A 是 Noether 环且 M 是理想 Hausdorff 的，则性质 (i)、(ii)、(iii)、(iv) 和 (v) 等价。

注。若 A/3 是一个域（这在应用中经常发生），则条件“ M/3M 是平坦的 (A/3)-模”对每个 A-模 M 都自动成立，从而简化了性质 (iii) 和 (iv) 的陈述；此外，在此情形下，性质 (v) 等价于：对于每个整数 n ≥ 1，M/3 \(^{n}\) M 是自由的 (A/3 \(^{n}\) )-模（第 II 章 § 3，第 2 节，命题 5 的推论 2）。

